\section{Аналитичность и топология}

\begin{axiom}[A9 · Дискретная аналитичность (Коши--Риман)]
\label{ax:A9}
$g$ --- оператор удержания дискретной голоморфности: голономия $\zeta=(\sum_N z)\cdot z^*$
вместо ложных скачков $\Arg(\sum z)-\Arg z$.
CR-единственность на связном $\carrierLattice$ задаёт глобальное поле (риманова поверхность $\carrierLattice\times S^1$).
\end{axiom}

\begin{axiom}[A10 · Топологическая квантованность заряда]
\label{ax:A10}
Устойчивый дефект --- полюс на $\partial(\hv)$. Топологический заряд
\[
  n=\frac{1}{2\pi}\oint_{\partial(\hv)} d\arg z \in \mathbb{Z}.
\]
\end{axiom}

\begin{axiom}[A11 · Устойчивость солитона]
\label{ax:A11}
Размазанное (не голоморфное) состояние не устойчиво: вращатель с $K_P$ и $\Delta\phi_{\min}$
выталкивает полюс --- самофокусировка vortex, не smear.
\end{axiom}

\begin{figure}[htbp]
  \centering
  \begin{tikzpicture}[carrier panel]
  \draw[carrier object] (0,0) rectangle (1,1);
  \foreach \x/\y in {0.15/0.15,0.85/0.15,0.85/0.85,0.15/0.85} {
    \draw[carrier arrow, carrier muted] (\x,\y) -- ++(0.12,0.12);
  }
\end{tikzpicture}

  \caption{$A_9$: дискретная голономия $\zeta$.}
  \label{fig:axiom-a9-holonomy}
\end{figure}

\begin{figure}[htbp]
  \centering
  \begin{tikzpicture}[carrier panel]
  \draw[carrier muted] (0,0) circle (0.9);
  \draw[carrier arrow] (0,0) -- (0.9,0);
  \draw[carrier object] (0.9,0) arc (0:300:0.9);
  \node[carrier site] at (0,0) {};
\end{tikzpicture}

  \caption{$A_{10}$: квантование намотки $\oint d\arg z=2\pi n$.}
  \label{fig:axiom-a10-winding}
\end{figure}

\begin{figure}[htbp]
  \centering
  \begin{tikzpicture}[carrier panel]
  \draw[carrier muted, dashed, domain=0.05:1.2, samples=50] plot (\x, {exp(-((\x-0.7)^2)/0.25)});
  \draw[carrier object, domain=0.05:1.2, samples=50] plot (\x, {exp(-(\x*\x)/0.03)});
  \node[carrier muted, font=\scriptsize] at (0.55,0.72) {размазано};
  \node[font=\scriptsize] at (0.15,0.88) {полюс};
  \node[carrier muted, font=\scriptsize] at (0.6,-0.4) {$r$};
\end{tikzpicture}

  \caption{$A_{11}$: устойчивый полюс (anti-smear), не размазанное состояние.}
  \label{fig:axiom-a11-vortex}
\end{figure}

